% Chapter: Proof Theory Foundations
% Part of the CCTT Monograph

\chapter{Proof Theory: Foundations}
\label{ch:proof-theory-foundations}

The preceding chapters established CCTT's ability to detect equivalences
via Z3 decision procedures and \v{C}ech cohomology.  But these techniques
have a fundamental limitation: they cannot handle \emph{algorithmic}
equivalences---cases where two programs compute the same function via
radically different strategies (e.g., dynamic programming vs.\ memoized
recursion, BFS vs.\ Union-Find, merge sort vs.\ quick sort).

In this chapter we develop a \textbf{proof theory} for CCTT: a formal
system in which a human (or LLM) can write an explicit proof term that
CCTT can \emph{mechanically verify}.  The proof checker is purely
structural---no sampling, no oracles, no heuristics.  If the proof
checks, the equivalence is \emph{proven}.

\section{Motivation: Why a Proof Theory?}

\begin{enumerate}
\item \textbf{Z3 is incomplete for algorithms.}
Z3 operates on quantifier-free formulas in decidable fragments.
It cannot reason about loop invariants, recursion depth, or
algorithmic structure directly.

\item \textbf{Bounded testing is not a proof.}
Even with millions of test inputs, a bounded test cannot guarantee
equivalence.  Adversarial inputs may exist outside the tested domain.

\item \textbf{Explicit evidence enables trust.}
A proof term is a \emph{certificate}---anyone can re-verify it.
Unlike a Z3 model or a sampling run, a proof term is permanent,
portable, and auditable.

\item \textbf{Compositionality.}
Proofs compose: if we prove $A \equiv B$ and $B \equiv C$, then
$\mathsf{Trans}$ gives us $A \equiv C$.  No ``leaps of faith.''
\end{enumerate}

\section{The Language of Proof Terms}

\begin{definition}[Proof Term]
A \textbf{proof term} is a syntactic witness that two OTerms are equal
in the CCTT equational theory.  We write
\[
  \Gamma \vdash p : a \equiv b
\]
where $\Gamma$ is a context of definitions and assumptions, $a$ and $b$
are OTerms, and $p$ is the proof term.
\end{definition}

\subsection{Structural Rules}

The structural rules form the core of the identity type.

\begin{definition}[Reflexivity]
$\mathsf{Refl}(t)$ proves $t \equiv t$.
\[
\frac{}{\Gamma \vdash \mathsf{Refl}(t) : t \equiv t} \quad (\text{refl})
\]
\end{definition}

\begin{definition}[Symmetry]
If $p : a \equiv b$, then $\mathsf{Sym}(p) : b \equiv a$.
\[
\frac{\Gamma \vdash p : a \equiv b}
     {\Gamma \vdash \mathsf{Sym}(p) : b \equiv a} \quad (\text{sym})
\]
\end{definition}

\begin{definition}[Transitivity]
If $p : a \equiv b$ and $q : b \equiv c$, then
$\mathsf{Trans}(p, q) : a \equiv c$.
\[
\frac{\Gamma \vdash p : a \equiv b \qquad \Gamma \vdash q : b \equiv c}
     {\Gamma \vdash \mathsf{Trans}(p, q) : a \equiv c} \quad (\text{trans})
\]
\end{definition}

\begin{definition}[Congruence]
If $p_i : a_i \equiv b_i$ for $i = 1, \ldots, n$, then
$\mathsf{Cong}(f, p_1, \ldots, p_n) : f(a_1,\ldots,a_n) \equiv f(b_1,\ldots,b_n)$.
\[
\frac{\Gamma \vdash p_i : a_i \equiv b_i \quad (i = 1, \ldots, n)}
     {\Gamma \vdash \mathsf{Cong}(f, \vec{p}) : f(\vec{a}) \equiv f(\vec{b})}
     \quad (\text{cong})
\]
\end{definition}

\subsection{Computation Rules}

\begin{definition}[$\beta$-reduction]
For a lambda $\lambda x.\, \mathit{body}$ applied to argument $\mathit{arg}$:
\[
\frac{}{\Gamma \vdash \mathsf{Beta}(\lambda x.\mathit{body}, \mathit{arg})
  : (\lambda x.\mathit{body})(\mathit{arg}) \equiv \mathit{body}[x := \mathit{arg}]}
  \quad (\beta)
\]
\end{definition}

\begin{definition}[$\delta$-reduction]
Unfold a named definition $f$:
\[
\frac{f(x_1, \ldots, x_n) \triangleq \mathit{body} \in \Gamma}
     {\Gamma \vdash \mathsf{Delta}(f, \vec{a})
       : f(\vec{a}) \equiv \mathit{body}[\vec{x} := \vec{a}]}
     \quad (\delta)
\]
\end{definition}

\begin{definition}[$\eta$-reduction]
When $x$ is not free in $f$:
\[
\frac{x \notin \mathrm{FV}(f)}
     {\Gamma \vdash \mathsf{Eta}(f)
       : \lambda x.\, f(x) \equiv f}
     \quad (\eta)
\]
\end{definition}

\subsection{Induction Principles}

\begin{definition}[Natural Number Induction]
To prove $\forall n.\, P(n)$:
\[
\frac{\Gamma \vdash p_0 : P(0) \qquad
      \Gamma, k : \Nat, \mathit{IH}: P(k) \vdash p_s : P(k+1)}
     {\Gamma \vdash \mathsf{NatInd}(p_0, p_s, n) : \forall n.\, P(n)}
     \quad (\Nat\text{-ind})
\]
\end{definition}

\begin{definition}[List Induction]
To prove $\forall \mathit{xs}.\, P(\mathit{xs})$:
\[
\frac{\Gamma \vdash p_\varepsilon : P([]) \qquad
      \Gamma, x, \mathit{xs}, \mathit{IH}: P(\mathit{xs}) \vdash p_c
        : P(x :: \mathit{xs})}
     {\Gamma \vdash \mathsf{ListInd}(p_\varepsilon, p_c, \mathit{xs})
       : \forall \mathit{xs}.\, P(\mathit{xs})}
     \quad (\mathsf{List}\text{-ind})
\]
\end{definition}

\begin{definition}[Well-Founded Induction]
Given a measure function $\mu$ into a well-ordered set:
\[
\frac{\Gamma, (\forall y.\, \mu(y) < \mu(x) \Rightarrow P(y)) \vdash p : P(x)}
     {\Gamma \vdash \mathsf{WFInd}(\mu, p) : \forall x.\, P(x)}
     \quad (\mathsf{WF}\text{-ind})
\]
\end{definition}

\subsection{Axiom Application}

\begin{definition}[Axiom Application]
Apply a named CCTT axiom $\mathcal{A}$ as a rewrite:
\[
\frac{\mathcal{A} : \ell \equiv r \quad
      \sigma : \mathrm{Var} \to \mathsf{OTerm}}
     {\Gamma \vdash \mathsf{AxiomApp}(\mathcal{A}, \sigma, t)
       : t[\ell\sigma] \equiv t[r\sigma]}
     \quad (\text{axiom})
\]
\end{definition}

\subsection{Loop and Recursion Reasoning}

\begin{definition}[Loop Invariant]
A Hoare-style loop proof with invariant $I$:
\[
\frac{p_{\mathit{init}} : I(s_0) \qquad
      p_{\mathit{pres}} : I(s) \wedge G(s) \Rightarrow I(\mathit{step}(s)) \qquad
      p_{\mathit{term}} : \mu \text{ decreases} \qquad
      p_{\mathit{post}} : I(s) \wedge \neg G(s) \Rightarrow Q(s)}
     {\Gamma \vdash \mathsf{LoopInv}(I, p_{\mathit{init}}, p_{\mathit{pres}},
       p_{\mathit{term}}, p_{\mathit{post}}) : \text{loop correct}}
\]
\end{definition}

\begin{definition}[Simulation]
Prove two programs equivalent by a simulation relation $R$:
\[
\frac{p_0 : R(s^f_0, s^g_0) \qquad
      p_s : R(s^f, s^g) \Rightarrow R(\mathit{step}_f(s^f), \mathit{step}_g(s^g)) \qquad
      p_o : R(s^f, s^g) \Rightarrow \mathit{out}_f(s^f) = \mathit{out}_g(s^g)}
     {\Gamma \vdash \mathsf{Sim}(R, p_0, p_s, p_o) : f \equiv g}
     \quad (\text{sim})
\]
\end{definition}

\begin{definition}[Abstraction Refinement]
Prove equivalence by showing both programs refine the same
abstract specification:
\[
\frac{p_f : \alpha \circ f = \mathit{spec} \qquad
      p_g : \alpha \circ g = \mathit{spec}}
     {\Gamma \vdash \mathsf{AbsRef}(\mathit{spec}, p_f, p_g) : f \equiv g}
     \quad (\text{abs-ref})
\]
\end{definition}

\subsection{Algebraic Reasoning}

\begin{definition}[Commutative Diagram]
If $h \circ f = g \circ k$ and $h, k$ are isomorphisms, then $f \equiv g$
up to encoding:
\[
\frac{p_c : h \circ f = g \circ k \qquad p_i : h, k \text{ are isos}}
     {\Gamma \vdash \mathsf{CommDiag}(p_c, p_i) : f \equiv g}
     \quad (\text{diag})
\]
\end{definition}

\subsection{Z3 Discharge}

\begin{definition}[Z3 Discharge]
For decidable subgoals in quantifier-free fragments:
\[
\frac{\varphi \in \mathrm{QF\_LIA} \cup \mathrm{QF\_NIA} \cup \cdots \qquad
      \mathrm{Z3}(\neg\varphi) = \mathsf{unsat}}
     {\Gamma \vdash \mathsf{Z3}(\varphi) : \varphi \text{ holds}}
     \quad (\text{z3})
\]
\textbf{Critical}: if Z3 returns $\mathsf{unknown}$ (timeout), the proof is
\textbf{rejected}---not ``probably true.''
\end{definition}

\subsection{Cohomological Reasoning}

\begin{definition}[Fiber Decomposition]
Prove equivalence by proving on each fiber separately:
\[
\frac{\forall \tau \in \mathrm{Fibers}.\;
      p_\tau : \restr{f}{\tau} \equiv \restr{g}{\tau}}
     {\Gamma \vdash \mathsf{FiberDecomp}(\{p_\tau\})
       : f \equiv g}
     \quad (\text{fiber})
\]
This is the sheaf condition: a global section is determined by its
restrictions to an open cover.
\end{definition}

\begin{definition}[\v{C}ech Gluing]
Glue local equivalences via \v{C}ech cohomology:
\[
\frac{\{p_i : \restr{f}{U_i} \equiv \restr{g}{U_i}\} \qquad
      \{p_{ij} : \text{compatible on } U_i \cap U_j\} \qquad
      \Hone = 0}
     {\Gamma \vdash \mathsf{CechGlue}(\vec{p}, \vec{p_{ij}})
       : f \equiv g}
     \quad (\text{glue})
\]
\end{definition}


\section{The Proof Checker}

The proof checker is a recursive function
$\mathsf{check} : \mathsf{ProofTerm} \times \mathsf{OTerm} \times \mathsf{OTerm}
\to \{\top, \bot\}$.

\begin{theorem}[Properties of the Checker]
The proof checker satisfies:
\begin{enumerate}
\item \textbf{Soundness}: If $\mathsf{check}(p, a, b) = \top$, then
  $\den{a} = \den{b}$ in the denotational semantics.
\item \textbf{Determinism}: The same proof term always produces the
  same result.
\item \textbf{Termination}: $\mathsf{check}$ always returns in bounded
  time (recursion depth $\leq 1000$; Z3 calls bounded by timeout).
\item \textbf{Non-sampling}: The checker never generates random inputs
  or performs bounded testing.
\end{enumerate}
\end{theorem}

\begin{proof}[Proof sketch]
Soundness follows from the correctness of each rule.
\begin{itemize}
\item $\mathsf{Refl}$: canonical form equality implies denotational equality.
\item $\mathsf{Sym}$, $\mathsf{Trans}$: the denotational equality relation is
  an equivalence relation.
\item $\mathsf{Cong}$: denotational semantics is compositional; equal
  arguments produce equal results.
\item $\beta, \delta, \eta$: reduction rules preserve denotational semantics
  by construction.
\item $\mathsf{NatInd}$, $\mathsf{ListInd}$: structural induction is sound
  for the natural numbers (Peano) and lists (free algebra).
\item $\mathsf{Z3Discharge}$: Z3 returning UNSAT for $\neg\varphi$ is a
  machine-checked proof of $\varphi$ in the given fragment.
\item $\mathsf{Simulation}$: a valid bisimulation relation implies
  observational equivalence.
\item $\mathsf{AbsRef}$: if $\alpha \circ f = \alpha \circ g$ and
  $\alpha$ is injective (or both sides equal the spec), then $f \equiv g$.
\end{itemize}

Determinism follows because the checker is purely functional (no
mutable state, no randomness).

Termination follows because:
(a)~proof depth is bounded by \texttt{MAX\_DEPTH = 1000};
(b)~Z3 calls have a timeout;
(c)~each recursive call strictly decreases the proof tree.
\end{proof}


\section{Soundness Theorem}

\begin{theorem}[Soundness of CCTT Proof Theory]
\label{thm:soundness}
If $\Gamma \vdash p : a \equiv b$ is derivable, then for all
environments $\rho$ satisfying $\Gamma$:
\[
  \den{a}_\rho = \den{b}_\rho
\]
\end{theorem}

\begin{proof}
By induction on the structure of the derivation $p$.
The base cases ($\mathsf{Refl}$, $\mathsf{Beta}$, $\mathsf{Delta}$,
$\mathsf{Eta}$, $\mathsf{Z3Discharge}$, $\mathsf{AxiomApp}$) hold
by the correctness of canonical-form computation, substitution,
$\eta$-law, Z3 as a decision procedure, and the soundness of each
CCTT axiom (established in the preceding chapters).

The inductive cases ($\mathsf{Sym}$, $\mathsf{Trans}$, $\mathsf{Cong}$,
$\mathsf{NatInd}$, $\mathsf{ListInd}$, $\mathsf{WFInd}$,
$\mathsf{Simulation}$, $\mathsf{AbsRef}$, $\mathsf{FiberDecomp}$,
$\mathsf{CechGlue}$) follow from the IH applied to sub-proofs, plus
the soundness of the corresponding mathematical principles (equational
reasoning, induction, bisimulation, refinement, sheaf descent).
\end{proof}


\section{Comparison with Other Proof Systems}

\subsection{Martin-L\"of Identity Types}

In Martin-L\"of type theory (MLTT), the identity type $\mathrm{Id}_A(a,b)$
has a single constructor $\mathsf{refl}$ and an elimination principle
$\mathsf{J}$.  CCTT's proof terms are richer: we include computation
rules ($\beta, \delta, \eta$), induction principles, and domain-specific
rules (simulation, loop invariants) that would require additional axioms
or libraries in MLTT.

\subsection{Cubical Path Types}

In cubical type theory, a path $p : \Path_A\, a\, b$ is a function
$p : \interval \to A$ with $p(0) = a$ and $p(1) = b$.  CCTT's path
search (Chapter~\ref{ch:path-search}) constructs paths using axiom
rewrites; the proof theory provides a language for \emph{certifying}
that a path exists.

The key difference: cubical paths are computational (they compute via
Kan operations); CCTT proof terms are \emph{verification certificates}
that the checker validates.

\subsection{Lean Proofs}

CCTT proof terms can be \emph{embedded} into Lean proofs:
\begin{itemize}
\item $\mathsf{Refl}$ maps to \texttt{rfl}
\item $\mathsf{Sym}$ maps to \texttt{Eq.symm}
\item $\mathsf{Trans}$ maps to \texttt{Eq.trans}
\item $\mathsf{Cong}$ maps to \texttt{congrArg}
\item $\mathsf{NatInd}$ maps to \texttt{Nat.rec}
\item $\mathsf{Z3Discharge}$ maps to \texttt{omega} or \texttt{decide}
\end{itemize}

\begin{theorem}[Embedding into Lean]
Every valid CCTT proof term can be translated into a Lean~4 proof term
that type-checks in the Lean kernel.
\end{theorem}

This embedding provides a second layer of assurance: even if the CCTT
proof checker has a bug, the Lean kernel provides an independent
verification.


\section{Proof Term Serialization}

Proof terms can be serialized to JSON for storage and exchange:
\begin{lstlisting}
{
  "tag": "NatInduction",
  "base_case": {"tag": "Refl", "term": {"type": "OLit", "value": 1}},
  "inductive_step": {"tag": "Assume", "label": "IH_n", ...},
  "variable": "n"
}
\end{lstlisting}
The serialization is complete: $\mathsf{from\_json}(\mathsf{to\_json}(p)) = p$
for every proof term $p$.

Proof documents include the goal (lhs, rhs), the proof term, and
metadata (name, description, strategy).  They can be stored alongside
source files as \texttt{*.proof.json} files.

\section{Integration with the CCTT Pipeline}

The proof theory integrates with the existing checker as
``Strategy~0: Explicit Proof Verification.''  When a proof file
exists alongside a program pair, the checker verifies the proof
\emph{instead of} using heuristics:

\begin{enumerate}
\item Try explicit proof verification (proof file or registered proof)
\item Try denotational OTerm equivalence
\item Try Z3 structural equality
\item Try semantic fingerprinting
\item Try per-fiber Z3 + \v{C}ech $\Hone$
\item Bounded testing (fallback)
\end{enumerate}

A verified proof gives confidence 1.0---no sampling or heuristics
involved.
